%=======================================================================
%  APPENDIX B -- COMPLETE RESULTS
%
%  The per-domain tables carry the full confusion matrix for every one of
%  the 27 configurations. Chapter 6 shows only AUC-ROC to stay readable;
%  this appendix is where a reader checks that nothing was selected to
%  flatter the argument. Every table here is GENERATED from the pipeline
%  CSVs by Thesis/tools/make_tables.py.
%=======================================================================
\chapter{Complete Results}
\label{app:full-results}

\section{Detection Metrics by Domain}

The complete confusion matrix, accuracy, precision, recall, $F_1$ and AUC-ROC for each of the 27 configurations are given in Chapter~\ref{ch:results}: Table~\ref{tab:detail-robotic} for robotics, Table~\ref{tab:detail-healthcare} for healthcare and Table~\ref{tab:detail-industrial} for industrial. They are not reproduced here, since the discussion of the divergence between $F_1$ and AUC-ROC in that chapter depends on having them alongside it.

\section{Bootstrap Confidence Intervals}

Chapter~\ref{ch:results} argues the uncertainty case from Figure~\ref{fig:forest}, which shows how far the 27 intervals overlap without asking the reader to compare bounds by eye. Table~\ref{tab:bootstrap} is the same estimates in numerical form, for when an exact bound is wanted. The analytical Hanley--McNeil interval is printed alongside each one: it is narrower almost everywhere, and the gap between the two columns is the cost of assuming the score series is independent when it is not.

% GENERATED by Thesis/tools/make_tables.py -- do not edit by hand.
% Every figure below is read from a CSV written by the pipelines.
\begin{table}[htbp]
\centering
\footnotesize
\setlength{\tabcolsep}{3pt}
\caption[Moving-block bootstrap confidence intervals on AUC-ROC, with the analytical Hanley--McNeil (HM) interval alongside for comparison]{Moving-block bootstrap confidence intervals on AUC-ROC, with the analytical Hanley--McNeil (HM) interval alongside for comparison, which is what the HM low and HM high columns report. Blocks preserve the serial correlation of the score series, which an i.i.d.\ resample would destroy and which would otherwise make every interval spuriously narrow. $L$ is the block length, derived per series as $\max(2, \operatorname{round}(n^{1/3}))$. The bold AUC in each algorithm-domain block is the best of the three regression variants on that data; no comparison is drawn between blocks, which score different observations.}
\label{tab:bootstrap}
\begin{tabular}{lllrrrrrrr}
\toprule
Algorithm & Domain & Variant & AUC $\uparrow$ & CI low & CI high & $n$ & $L$ & HM low & HM high \\
\midrule
CAM-UV & Healthcare & Linear & \textbf{0.7861} & 0.7319 & 0.8325 & 3115 & 15 & 0.7703 & 0.8019 \\
CAM-UV & Healthcare & Poly-3 & 0.7691 & 0.7103 & 0.8196 & 3115 & 15 & 0.7528 & 0.7854 \\
CAM-UV & Healthcare & RBF & 0.7692 & 0.7161 & 0.8198 & 3115 & 15 & 0.7529 & 0.7855 \\
CAM-UV & Industrial & Linear & 0.8428 & 0.7601 & 0.9007 & 1323 & 11 & 0.8070 & 0.8787 \\
CAM-UV & Industrial & Poly-3 & 0.8212 & 0.7000 & 0.9153 & 1323 & 11 & 0.7816 & 0.8607 \\
CAM-UV & Industrial & RBF & \textbf{0.8922} & 0.8109 & 0.9416 & 1323 & 11 & 0.8656 & 0.9189 \\
CAM-UV & Robotics & Linear & 0.6497 & 0.5475 & 0.7383 & 1378 & 11 & 0.6209 & 0.6785 \\
CAM-UV & Robotics & Poly-3 & 0.6888 & 0.5952 & 0.7714 & 1378 & 11 & 0.6612 & 0.7164 \\
CAM-UV & Robotics & RBF & \textbf{0.7818} & 0.6915 & 0.8564 & 1378 & 11 & 0.7581 & 0.8055 \\
GES & Healthcare & Linear & \textbf{0.8114} & 0.7628 & 0.8540 & 3115 & 15 & 0.7965 & 0.8263 \\
GES & Healthcare & Poly-3 & 0.7867 & 0.7314 & 0.8379 & 3115 & 15 & 0.7709 & 0.8025 \\
GES & Healthcare & RBF & 0.7620 & 0.7061 & 0.8141 & 3115 & 15 & 0.7455 & 0.7785 \\
GES & Industrial & Linear & 0.8407 & 0.7463 & 0.9022 & 1323 & 11 & 0.8045 & 0.8769 \\
GES & Industrial & Poly-3 & 0.8849 & 0.8048 & 0.9355 & 1323 & 11 & 0.8568 & 0.9130 \\
GES & Industrial & RBF & \textbf{0.8994} & 0.8206 & 0.9508 & 1323 & 11 & 0.8741 & 0.9246 \\
GES & Robotics & Linear & 0.5906 & 0.4873 & 0.6761 & 1378 & 11 & 0.5604 & 0.6207 \\
GES & Robotics & Poly-3 & \textbf{0.6999} & 0.6017 & 0.7804 & 1378 & 11 & 0.6727 & 0.7271 \\
GES & Robotics & RBF & 0.6516 & 0.5472 & 0.7400 & 1378 & 11 & 0.6228 & 0.6803 \\
PCMCI & Healthcare & Linear & 0.6181 & 0.5536 & 0.6755 & 4806 & 17 & 0.6023 & 0.6339 \\
PCMCI & Healthcare & Poly-3 & 0.5075 & 0.4429 & 0.5663 & 4806 & 17 & 0.4909 & 0.5242 \\
PCMCI & Healthcare & RBF & \textbf{0.7147} & 0.6597 & 0.7640 & 4806 & 17 & 0.7004 & 0.7289 \\
PCMCI & Industrial & Linear & 0.8860 & 0.8537 & 0.9124 & 16716 & 26 & 0.8782 & 0.8939 \\
PCMCI & Industrial & Poly-3 & \textbf{0.9151} & 0.8880 & 0.9394 & 16716 & 26 & 0.9089 & 0.9213 \\
PCMCI & Industrial & RBF & 0.8495 & 0.8118 & 0.8796 & 16716 & 26 & 0.8398 & 0.8593 \\
PCMCI & Robotics & Linear & 0.7825 & 0.6321 & 0.8948 & 181 & 6 & 0.7170 & 0.8480 \\
PCMCI & Robotics & Poly-3 & 0.8053 & 0.6516 & 0.9233 & 181 & 6 & 0.7432 & 0.8674 \\
PCMCI & Robotics & RBF & \textbf{0.8752} & 0.7559 & 0.9587 & 181 & 6 & 0.8258 & 0.9246 \\
\bottomrule
\end{tabular}
\end{table}


\section{Complete Attribution Rankings}

Chapter~\ref{ch:results} states the robotics attribution result in a sentence and then leaves the per-fault rankings to this appendix. They are given twice over: first as the top-ranked variable for each of the three robotics attack classes, which is what the documented requirement is stated over, and then as the full top-three ranking for every attack class in every configuration and domain, so that the robotics result can be read against the whole set rather than in isolation. The ranking rule differs by domain and is defined in Section~\ref{sec:method-attribution}.

\subsection{Top-ranked variable per robotics attack}

Table~\ref{tab:attributions} ranks by the $L_2$ norm of the attack-window coefficient drift, the rule retained on robotics for parity with the reference implementation; the healthcare and industrial pipelines normalise by each variable's fault-free baseline instead.

% GENERATED by Thesis/tools/make_tables.py -- do not edit by hand.
% Every figure below is read from a CSV written by the pipelines.
\begin{table}[htbp]
\centering
\small
\caption[Top-ranked attributed variable for each robotics attack class, across all nine robotics configurations]{Top-ranked attributed variable for each robotics attack class, across all nine robotics configurations. The documented explainability requirement for this dataset is that LED attacks attribute to a tablet sensor.}
\label{tab:attributions}
\begin{tabular}{lllll}
\toprule
Algorithm & Variant & LedsControl & JointControl & WheelsControl \\
\midrule
\multirow{3}{*}{PCMCI} & Linear & TabletTouchX & KneePicthTemp & WheelFRTemp \\
 & Poly-3 & TabletTouchX & WheelBStiff & HeadPicthTemp \\
 & \textbf{RBF} & \textbf{TabletTouchX} & \textbf{KneePicthTemp} & \textbf{WheelBTemp} \\
\midrule
\multirow{3}{*}{GES} & Linear & TabletTouchX & KneePicthTemp & LElbowYawTemp \\
 & \textbf{Poly-3} & \textbf{TabletTouchX} & \textbf{KneePicthTemp} & \textbf{LShoulderRollTemp} \\
 & RBF & TabletTouchX & KneePicthTemp & LaserShovel03X \\
\midrule
\multirow{3}{*}{CAM-UV} & Linear & TabletTouchX & KneePicthTemp & WheelFRTemp \\
 & Poly-3 & TabletTouchX & KneePicthTemp & WheelFRStiff \\
 & \textbf{RBF} & \textbf{TabletTouchX} & \textbf{KneePicthTemp} & \textbf{KneePicthTemp} \\
\bottomrule
\end{tabular}
\end{table}


The attribution isolates the downstream tablet sensors for the \texttt{LedsControl} attack in 9 of 9 evaluated configurations, and recovers the thermal root cause of the mechanical \texttt{JointControl} attack in 8 of 9. The exception is the PCMCI polynomial variant, which attributes that anomaly to \texttt{WheelBStiff}, a wheel stiffness reading unrelated to the joint subsystem. This consensus is reached despite the algorithms generating entirely different graph structures and operating on different input dimensions.

Why the attribution lands on the tablet and not on the LEDs deserves setting out, since this is
the clearest case of the mechanism doing what it was built to do. For PCMCI the graph places the LED channels as parents of the tablet sensors. When the LED attack fires, the relationship predicting the tablet from the LEDs breaks, and the variable that becomes unpredictable is the child; the pipeline surfaces the downstream victim. For GES and CAM-UV the 35-variable variance cap excludes every LED channel, so those graphs cannot represent an LED-to-tablet relationship at all, and the tablet sensors still carry the largest attack-window coefficient drift. The ranking is not an artefact of the metric either: ordering by raw $L_2$ norm and by baseline-normalised RMS ratio both return \texttt{TabletTouchX} and \texttt{TabletTouchY} as the top two for this attack.

That agreement does not extend to the other two faults, and the divergence justifies the per-domain choice of metric declared in Section~\ref{sec:method-attribution}. On the \texttt{JointControl} anomaly the raw $L_2$ norm returns \texttt{KneePicthTemp}, the joint temperature the documented requirement names, whereas the baseline-normalised ratio returns \texttt{LedEarR3} and \texttt{LedEarL3}. The reason is scale instead of physiology: those two channels carry the twelfth and thirteenth smallest fault-free baselines of 244 variables, at 0.0268, against \texttt{KneePicthTemp} at rank 238 and 7.5380. Dividing by a baseline that small confers a 281-fold advantage on two quiet sensors and displaces the correct attribution entirely.

Neither rule dominates. On the \texttt{WheelsControl} attack the normalised ratio returns \texttt{WheelFRSpeed} and \texttt{WheelBSpeed}, which are arguably more apt than the raw norm's \texttt{WheelFRTemp} and \texttt{HeadPicthTemp}. The choice of attribution metric is load-bearing rather than cosmetic, and has to be made against the variance profile of the domain instead of adopted uniformly.

\subsection{Robotics}

Every robotics configuration against all three attack classes. The ranking is the raw $L_2$ norm of the attack-window drift, the rule retained on this domain for parity with the reference implementation.

% GENERATED by Thesis/tools/make_tables.py -- do not edit by hand.
% Every figure below is read from a CSV written by the pipelines.
\begin{table}[htbp]
\centering
\scriptsize
\setlength{\tabcolsep}{4pt}
\caption[Complete top-three attribution for the robotics domain, every configuration and every attack class]{Complete top-three attribution for the robotics domain, every configuration and every attack class. The ranking rule is stated in Section~\ref{sec:method-attribution}.}
\label{app:attr-robotic}
\begin{tabular}{llllll}
\toprule
Algorithm & Variant & Fault & First & Second & Third \\
\midrule
PCMCI & Linear & WheelsControl & \texttt{WheelFRTemp} & \texttt{HeadPicthTemp} & \texttt{RElbowYawTemp} \\
 & Linear & JointControl & \texttt{KneePicthTemp} & \texttt{WheelBTemp} & \texttt{WheelFRTemp} \\
 & Linear & LedsControl & \texttt{TabletTouchX} & \texttt{TabletTouchY} & \texttt{HeadPicthTemp} \\
 & Poly-3 & WheelsControl & \texttt{HeadPicthTemp} & \texttt{WheelBStiff} & \texttt{WheelBTemp} \\
 & Poly-3 & JointControl & \texttt{WheelBStiff} & \texttt{WheelFRStiff} & \texttt{WheelBTemp} \\
 & Poly-3 & LedsControl & \texttt{TabletTouchX} & \texttt{TabletTouchY} & \texttt{WheelBStiff} \\
 & RBF & WheelsControl & \texttt{WheelBTemp} & \texttt{WheelFRTemp} & \texttt{WheelBStiff} \\
 & RBF & JointControl & \texttt{KneePicthTemp} & \texttt{HipPicthTemp} & \texttt{WheelBTemp} \\
 & RBF & LedsControl & \texttt{TabletTouchX} & \texttt{TabletTouchY} & \texttt{HeadPicthTemp} \\
GES & Linear & WheelsControl & \texttt{LElbowYawTemp} & \texttt{WheelFLTemp} & \texttt{KneePicthTemp} \\
 & Linear & JointControl & \texttt{KneePicthTemp} & \texttt{HeadPicthTemp} & \texttt{LShoulderPitchTemp} \\
 & Linear & LedsControl & \texttt{TabletTouchX} & \texttt{TabletTouchY} & \texttt{HeadPicthTemp} \\
 & Poly-3 & WheelsControl & \texttt{LShoulderRollTemp} & \texttt{LaserR09X} & \texttt{RShoulderRollTemp} \\
 & Poly-3 & JointControl & \texttt{KneePicthTemp} & \texttt{HeadPicthTemp} & \texttt{LaserR07X} \\
 & Poly-3 & LedsControl & \texttt{TabletTouchX} & \texttt{TabletTouchY} & \texttt{HeadPicthTemp} \\
 & RBF & WheelsControl & \texttt{LaserShovel03X} & \texttt{LaserShovel01X} & \texttt{HeadPicthTemp} \\
 & RBF & JointControl & \texttt{KneePicthTemp} & \texttt{LaserR15X} & \texttt{LaserR12X} \\
 & RBF & LedsControl & \texttt{TabletTouchX} & \texttt{TabletTouchY} & \texttt{HeadPicthTemp} \\
CAM-UV & Linear & WheelsControl & \texttt{WheelFRTemp} & \texttt{WheelBTemp} & \texttt{WheelFRStiff} \\
 & Linear & JointControl & \texttt{KneePicthTemp} & \texttt{WheelFRStiff} & \texttt{WheelFRTemp} \\
 & Linear & LedsControl & \texttt{TabletTouchX} & \texttt{TabletTouchY} & \texttt{WheelBStiff} \\
 & Poly-3 & WheelsControl & \texttt{WheelFRStiff} & \texttt{KneePicthTemp} & \texttt{WheelFRTemp} \\
 & Poly-3 & JointControl & \texttt{KneePicthTemp} & \texttt{WheelFRTemp} & \texttt{WheelBTemp} \\
 & Poly-3 & LedsControl & \texttt{TabletTouchX} & \texttt{TabletTouchY} & \texttt{HeadPicthTemp} \\
 & RBF & WheelsControl & \texttt{KneePicthTemp} & \texttt{HipRollTemp} & \texttt{WheelFRTemp} \\
 & RBF & JointControl & \texttt{KneePicthTemp} & \texttt{HeadPicthTemp} & \texttt{WheelBTemp} \\
 & RBF & LedsControl & \texttt{TabletTouchX} & \texttt{TabletTouchY} & \texttt{KneePicthTemp} \\
\bottomrule
\end{tabular}
\end{table}


\subsection{Healthcare}

Every healthcare configuration against all twenty attack classes, ranked by the baseline-normalised RMS ratio rather than the raw norm.

% GENERATED by Thesis/tools/make_tables.py -- do not edit by hand.
% Every figure below is read from a CSV written by the pipelines.
{\scriptsize
\setlength{\tabcolsep}{4pt}
\begin{longtable}{llllll}
\caption[Complete top-three attribution for the healthcare domain, every configuration and every attack class]{Complete top-three attribution for the healthcare domain, every configuration and every attack class. The ranking rule is stated in Section~\ref{sec:method-attribution}.}\label{app:attr-healthcare}\\
\toprule
Algorithm & Variant & Fault & First & Second & Third \\
\midrule
\endfirsthead
\toprule
Algorithm & Variant & Fault & First & Second & Third \\
\midrule
\endhead
PCMCI & Linear & Attack 1 & \texttt{II} & \texttt{V6} & \texttt{V5} \\
 & Linear & Attack 2 & \texttt{V4} & \texttt{V3} & \texttt{V5} \\
 & Linear & Attack 3 & \texttt{V3} & \texttt{AVL} & \texttt{V6} \\
 & Linear & Attack 4 & \texttt{V3} & \texttt{AVL} & \texttt{V2} \\
 & Linear & Attack 5 & \texttt{AVR} & \texttt{V1} & \texttt{AVL} \\
 & Linear & Attack 6 & \texttt{V3} & \texttt{AVL} & \texttt{V2} \\
 & Linear & Attack 7 & \texttt{V3} & \texttt{V2} & \texttt{V1} \\
 & Linear & Attack 8 & \texttt{V4} & \texttt{V3} & \texttt{AVL} \\
 & Linear & Attack 9 & \texttt{V6} & \texttt{V5} & \texttt{V3} \\
 & Linear & Attack 10 & \texttt{V1} & \texttt{V2} & \texttt{V3} \\
 & Linear & Attack 11 & \texttt{AVL} & \texttt{AVR} & \texttt{V3} \\
 & Linear & Attack 12 & \texttt{V3} & \texttt{AVL} & \texttt{V5} \\
 & Linear & Attack 13 & \texttt{II} & \texttt{AVF} & \texttt{V4} \\
 & Linear & Attack 14 & \texttt{V5} & \texttt{V4} & \texttt{AVL} \\
 & Linear & Attack 15 & \texttt{V6} & \texttt{AVR} & \texttt{AVL} \\
 & Linear & Attack 16 & \texttt{II} & \texttt{V3} & \texttt{III} \\
 & Linear & Attack 17 & \texttt{V5} & \texttt{V4} & \texttt{II} \\
 & Linear & Attack 18 & \texttt{V4} & \texttt{V3} & \texttt{V5} \\
 & Linear & Attack 19 & \texttt{AVL} & \texttt{V1} & \texttt{I} \\
 & Linear & Attack 20 & \texttt{V6} & \texttt{V5} & \texttt{AVR} \\
 & Poly-3 & Attack 1 & \texttt{AVF} & \texttt{II} & \texttt{AVL} \\
 & Poly-3 & Attack 2 & \texttt{V4} & \texttt{V3} & \texttt{V5} \\
 & Poly-3 & Attack 3 & \texttt{V5} & \texttt{V3} & \texttt{V4} \\
 & Poly-3 & Attack 4 & \texttt{AVF} & \texttt{II} & \texttt{AVL} \\
 & Poly-3 & Attack 5 & \texttt{AVR} & \texttt{V1} & \texttt{I} \\
 & Poly-3 & Attack 6 & \texttt{V3} & \texttt{V4} & \texttt{V5} \\
 & Poly-3 & Attack 7 & \texttt{V3} & \texttt{V4} & \texttt{V5} \\
 & Poly-3 & Attack 8 & \texttt{V4} & \texttt{V6} & \texttt{V3} \\
 & Poly-3 & Attack 9 & \texttt{V3} & \texttt{V2} & \texttt{III} \\
 & Poly-3 & Attack 10 & \texttt{V1} & \texttt{AVF} & \texttt{V2} \\
 & Poly-3 & Attack 11 & \texttt{AVR} & \texttt{V3} & \texttt{V2} \\
 & Poly-3 & Attack 12 & \texttt{V5} & \texttt{V6} & \texttt{AVL} \\
 & Poly-3 & Attack 13 & \texttt{AVF} & \texttt{V2} & \texttt{II} \\
 & Poly-3 & Attack 14 & \texttt{AVR} & \texttt{V5} & \texttt{II} \\
 & Poly-3 & Attack 15 & \texttt{V6} & \texttt{II} & \texttt{I} \\
 & Poly-3 & Attack 16 & \texttt{AVR} & \texttt{II} & \texttt{V3} \\
 & Poly-3 & Attack 17 & \texttt{V5} & \texttt{V6} & \texttt{II} \\
 & Poly-3 & Attack 18 & \texttt{V3} & \texttt{V4} & \texttt{V2} \\
 & Poly-3 & Attack 19 & \texttt{II} & \texttt{III} & \texttt{V4} \\
 & Poly-3 & Attack 20 & \texttt{V5} & \texttt{V6} & \texttt{V3} \\
 & RBF & Attack 1 & \texttt{III} & \texttt{AVF} & \texttt{V6} \\
 & RBF & Attack 2 & \texttt{V2} & \texttt{I} & \texttt{V3} \\
 & RBF & Attack 3 & \texttt{V2} & \texttt{V6} & \texttt{V5} \\
 & RBF & Attack 4 & \texttt{V2} & \texttt{I} & \texttt{AVL} \\
 & RBF & Attack 5 & \texttt{V2} & \texttt{V1} & \texttt{I} \\
 & RBF & Attack 6 & \texttt{V2} & \texttt{I} & \texttt{V3} \\
 & RBF & Attack 7 & \texttt{I} & \texttt{AVR} & \texttt{V2} \\
 & RBF & Attack 8 & \texttt{V2} & \texttt{V3} & \texttt{I} \\
 & RBF & Attack 9 & \texttt{V6} & \texttt{V5} & \texttt{V3} \\
 & RBF & Attack 10 & \texttt{V1} & \texttt{V3} & \texttt{V2} \\
 & RBF & Attack 11 & \texttt{I} & \texttt{V2} & \texttt{AVL} \\
 & RBF & Attack 12 & \texttt{I} & \texttt{AVL} & \texttt{V2} \\
 & RBF & Attack 13 & \texttt{V2} & \texttt{III} & \texttt{V3} \\
 & RBF & Attack 14 & \texttt{I} & \texttt{V5} & \texttt{AVR} \\
 & RBF & Attack 15 & \texttt{V6} & \texttt{AVR} & \texttt{I} \\
 & RBF & Attack 16 & \texttt{V3} & \texttt{III} & \texttt{AVF} \\
 & RBF & Attack 17 & \texttt{AVL} & \texttt{III} & \texttt{I} \\
 & RBF & Attack 18 & \texttt{V4} & \texttt{V5} & \texttt{V3} \\
 & RBF & Attack 19 & \texttt{I} & \texttt{V1} & \texttt{AVL} \\
 & RBF & Attack 20 & \texttt{V6} & \texttt{V5} & \texttt{I} \\
GES & Linear & Attack 1 & \texttt{AVL} & \texttt{AVF} & \texttt{III} \\
 & Linear & Attack 2 & \texttt{V4} & \texttt{II} & \texttt{AVF} \\
 & Linear & Attack 3 & \texttt{V4} & \texttt{V3} & \texttt{III} \\
 & Linear & Attack 4 & \texttt{V3} & \texttt{AVF} & \texttt{III} \\
 & Linear & Attack 5 & \texttt{AVL} & \texttt{I} & \texttt{V6} \\
 & Linear & Attack 6 & \texttt{V4} & \texttt{AVF} & \texttt{III} \\
 & Linear & Attack 7 & \texttt{I} & \texttt{II} & \texttt{V4} \\
 & Linear & Attack 8 & \texttt{V4} & \texttt{V3} & \texttt{V1} \\
 & Linear & Attack 9 & \texttt{V3} & \texttt{II} & \texttt{AVF} \\
 & Linear & Attack 10 & \texttt{AVF} & \texttt{V4} & \texttt{V3} \\
 & Linear & Attack 11 & \texttt{III} & \texttt{AVF} & \texttt{V4} \\
 & Linear & Attack 12 & \texttt{I} & \texttt{AVL} & \texttt{V6} \\
 & Linear & Attack 13 & \texttt{III} & \texttt{AVF} & \texttt{II} \\
 & Linear & Attack 14 & \texttt{II} & \texttt{III} & \texttt{AVF} \\
 & Linear & Attack 15 & \texttt{V6} & \texttt{V3} & \texttt{AVF} \\
 & Linear & Attack 16 & \texttt{V4} & \texttt{V3} & \texttt{AVF} \\
 & Linear & Attack 17 & \texttt{AVL} & \texttt{AVF} & \texttt{V6} \\
 & Linear & Attack 18 & \texttt{V6} & \texttt{V4} & \texttt{AVF} \\
 & Linear & Attack 19 & \texttt{I} & \texttt{III} & \texttt{V6} \\
 & Linear & Attack 20 & \texttt{V3} & \texttt{V4} & \texttt{V6} \\
 & Poly-3 & Attack 1 & \texttt{AVL} & \texttt{V6} & \texttt{I} \\
 & Poly-3 & Attack 2 & \texttt{V4} & \texttt{II} & \texttt{AVF} \\
 & Poly-3 & Attack 3 & \texttt{V3} & \texttt{III} & \texttt{AVF} \\
 & Poly-3 & Attack 4 & \texttt{II} & \texttt{AVF} & \texttt{III} \\
 & Poly-3 & Attack 5 & \texttt{III} & \texttt{V1} & \texttt{AVL} \\
 & Poly-3 & Attack 6 & \texttt{AVF} & \texttt{III} & \texttt{II} \\
 & Poly-3 & Attack 7 & \texttt{I} & \texttt{V3} & \texttt{II} \\
 & Poly-3 & Attack 8 & \texttt{I} & \texttt{AVL} & \texttt{V4} \\
 & Poly-3 & Attack 9 & \texttt{V3} & \texttt{V6} & \texttt{V4} \\
 & Poly-3 & Attack 10 & \texttt{V1} & \texttt{V3} & \texttt{AVL} \\
 & Poly-3 & Attack 11 & \texttt{III} & \texttt{AVF} & \texttt{II} \\
 & Poly-3 & Attack 12 & \texttt{I} & \texttt{AVL} & \texttt{V6} \\
 & Poly-3 & Attack 13 & \texttt{AVF} & \texttt{V3} & \texttt{III} \\
 & Poly-3 & Attack 14 & \texttt{I} & \texttt{II} & \texttt{AVF} \\
 & Poly-3 & Attack 15 & \texttt{V6} & \texttt{I} & \texttt{V3} \\
 & Poly-3 & Attack 16 & \texttt{V3} & \texttt{V4} & \texttt{AVF} \\
 & Poly-3 & Attack 17 & \texttt{III} & \texttt{AVL} & \texttt{I} \\
 & Poly-3 & Attack 18 & \texttt{V4} & \texttt{V3} & \texttt{V6} \\
 & Poly-3 & Attack 19 & \texttt{I} & \texttt{III} & \texttt{V1} \\
 & Poly-3 & Attack 20 & \texttt{V3} & \texttt{V4} & \texttt{V6} \\
 & RBF & Attack 1 & \texttt{III} & \texttt{II} & \texttt{AVF} \\
 & RBF & Attack 2 & \texttt{V4} & \texttt{AVL} & \texttt{I} \\
 & RBF & Attack 3 & \texttt{V3} & \texttt{V4} & \texttt{V6} \\
 & RBF & Attack 4 & \texttt{V6} & \texttt{V3} & \texttt{V4} \\
 & RBF & Attack 5 & \texttt{AVL} & \texttt{I} & \texttt{V4} \\
 & RBF & Attack 6 & \texttt{V6} & \texttt{V4} & \texttt{V3} \\
 & RBF & Attack 7 & \texttt{I} & \texttt{II} & \texttt{V3} \\
 & RBF & Attack 8 & \texttt{V4} & \texttt{I} & \texttt{V1} \\
 & RBF & Attack 9 & \texttt{III} & \texttt{V3} & \texttt{AVF} \\
 & RBF & Attack 10 & \texttt{V3} & \texttt{AVF} & \texttt{AVL} \\
 & RBF & Attack 11 & \texttt{V3} & \texttt{I} & \texttt{V4} \\
 & RBF & Attack 12 & \texttt{I} & \texttt{AVL} & \texttt{V6} \\
 & RBF & Attack 13 & \texttt{AVL} & \texttt{V3} & \texttt{III} \\
 & RBF & Attack 14 & \texttt{AVF} & \texttt{I} & \texttt{AVL} \\
 & RBF & Attack 15 & \texttt{V6} & \texttt{I} & \texttt{AVL} \\
 & RBF & Attack 16 & \texttt{V3} & \texttt{V4} & \texttt{AVL} \\
 & RBF & Attack 17 & \texttt{III} & \texttt{AVF} & \texttt{II} \\
 & RBF & Attack 18 & \texttt{V6} & \texttt{V4} & \texttt{I} \\
 & RBF & Attack 19 & \texttt{I} & \texttt{V6} & \texttt{AVL} \\
 & RBF & Attack 20 & \texttt{V3} & \texttt{V4} & \texttt{V6} \\
CAM-UV & Linear & Attack 1 & \texttt{I} & \texttt{AVF} & \texttt{AVL} \\
 & Linear & Attack 2 & \texttt{V4} & \texttt{V5} & \texttt{AVR} \\
 & Linear & Attack 3 & \texttt{V3} & \texttt{V2} & \texttt{III} \\
 & Linear & Attack 4 & \texttt{V2} & \texttt{I} & \texttt{V3} \\
 & Linear & Attack 5 & \texttt{V5} & \texttt{V4} & \texttt{V2} \\
 & Linear & Attack 6 & \texttt{V6} & \texttt{AVF} & \texttt{I} \\
 & Linear & Attack 7 & \texttt{AVR} & \texttt{I} & \texttt{AVL} \\
 & Linear & Attack 8 & \texttt{V4} & \texttt{V3} & \texttt{V2} \\
 & Linear & Attack 9 & \texttt{V3} & \texttt{V2} & \texttt{AVF} \\
 & Linear & Attack 10 & \texttt{V3} & \texttt{V2} & \texttt{V6} \\
 & Linear & Attack 11 & \texttt{AVL} & \texttt{I} & \texttt{V6} \\
 & Linear & Attack 12 & \texttt{I} & \texttt{V5} & \texttt{V4} \\
 & Linear & Attack 13 & \texttt{AVL} & \texttt{III} & \texttt{AVF} \\
 & Linear & Attack 14 & \texttt{II} & \texttt{AVR} & \texttt{AVF} \\
 & Linear & Attack 15 & \texttt{V6} & \texttt{I} & \texttt{AVR} \\
 & Linear & Attack 16 & \texttt{V2} & \texttt{V3} & \texttt{I} \\
 & Linear & Attack 17 & \texttt{V5} & \texttt{AVL} & \texttt{V6} \\
 & Linear & Attack 18 & \texttt{V5} & \texttt{V6} & \texttt{I} \\
 & Linear & Attack 19 & \texttt{AVL} & \texttt{I} & \texttt{III} \\
 & Linear & Attack 20 & \texttt{V3} & \texttt{V2} & \texttt{V4} \\
 & Poly-3 & Attack 1 & \texttt{AVL} & \texttt{I} & \texttt{V6} \\
 & Poly-3 & Attack 2 & \texttt{V4} & \texttt{II} & \texttt{AVR} \\
 & Poly-3 & Attack 3 & \texttt{V6} & \texttt{III} & \texttt{V2} \\
 & Poly-3 & Attack 4 & \texttt{I} & \texttt{AVF} & \texttt{III} \\
 & Poly-3 & Attack 5 & \texttt{V2} & \texttt{AVL} & \texttt{V1} \\
 & Poly-3 & Attack 6 & \texttt{AVF} & \texttt{V6} & \texttt{III} \\
 & Poly-3 & Attack 7 & \texttt{I} & \texttt{AVR} & \texttt{II} \\
 & Poly-3 & Attack 8 & \texttt{AVL} & \texttt{V2} & \texttt{V4} \\
 & Poly-3 & Attack 9 & \texttt{V6} & \texttt{V3} & \texttt{AVF} \\
 & Poly-3 & Attack 10 & \texttt{V3} & \texttt{I} & \texttt{V2} \\
 & Poly-3 & Attack 11 & \texttt{I} & \texttt{AVF} & \texttt{AVL} \\
 & Poly-3 & Attack 12 & \texttt{I} & \texttt{AVL} & \texttt{V5} \\
 & Poly-3 & Attack 13 & \texttt{III} & \texttt{AVF} & \texttt{AVL} \\
 & Poly-3 & Attack 14 & \texttt{AVR} & \texttt{II} & \texttt{V5} \\
 & Poly-3 & Attack 15 & \texttt{V6} & \texttt{AVR} & \texttt{II} \\
 & Poly-3 & Attack 16 & \texttt{V3} & \texttt{V2} & \texttt{V4} \\
 & Poly-3 & Attack 17 & \texttt{AVL} & \texttt{III} & \texttt{V5} \\
 & Poly-3 & Attack 18 & \texttt{V5} & \texttt{V4} & \texttt{V6} \\
 & Poly-3 & Attack 19 & \texttt{AVL} & \texttt{III} & \texttt{AVF} \\
 & Poly-3 & Attack 20 & \texttt{V3} & \texttt{V4} & \texttt{V2} \\
 & RBF & Attack 1 & \texttt{III} & \texttt{II} & \texttt{AVF} \\
 & RBF & Attack 2 & \texttt{V4} & \texttt{V5} & \texttt{AVL} \\
 & RBF & Attack 3 & \texttt{V2} & \texttt{V3} & \texttt{V5} \\
 & RBF & Attack 4 & \texttt{V5} & \texttt{V6} & \texttt{V4} \\
 & RBF & Attack 5 & \texttt{V2} & \texttt{I} & \texttt{V5} \\
 & RBF & Attack 6 & \texttt{V6} & \texttt{V5} & \texttt{V4} \\
 & RBF & Attack 7 & \texttt{I} & \texttt{AVR} & \texttt{AVL} \\
 & RBF & Attack 8 & \texttt{V4} & \texttt{V2} & \texttt{V1} \\
 & RBF & Attack 9 & \texttt{II} & \texttt{III} & \texttt{AVF} \\
 & RBF & Attack 10 & \texttt{V3} & \texttt{V2} & \texttt{AVF} \\
 & RBF & Attack 11 & \texttt{V3} & \texttt{AVL} & \texttt{I} \\
 & RBF & Attack 12 & \texttt{I} & \texttt{AVL} & \texttt{AVR} \\
 & RBF & Attack 13 & \texttt{III} & \texttt{II} & \texttt{AVL} \\
 & RBF & Attack 14 & \texttt{AVF} & \texttt{I} & \texttt{AVR} \\
 & RBF & Attack 15 & \texttt{V6} & \texttt{I} & \texttt{V2} \\
 & RBF & Attack 16 & \texttt{V2} & \texttt{V3} & \texttt{V4} \\
 & RBF & Attack 17 & \texttt{V5} & \texttt{III} & \texttt{AVL} \\
 & RBF & Attack 18 & \texttt{V5} & \texttt{V6} & \texttt{V4} \\
 & RBF & Attack 19 & \texttt{I} & \texttt{V6} & \texttt{V5} \\
 & RBF & Attack 20 & \texttt{V5} & \texttt{V3} & \texttt{V4} \\
\bottomrule
\end{longtable}}


\subsection{Industrial}

Every industrial configuration against all twenty fault classes, under the same baseline-normalised rule the healthcare tables use.

% GENERATED by Thesis/tools/make_tables.py -- do not edit by hand.
% Every figure below is read from a CSV written by the pipelines.
{\scriptsize
\setlength{\tabcolsep}{4pt}
\begin{longtable}{llllll}
\caption[Complete top-three attribution for the industrial domain, every configuration and every attack class]{Complete top-three attribution for the industrial domain, every configuration and every attack class. The ranking rule is stated in Section~\ref{sec:method-attribution}.}\label{app:attr-industrial}\\
\toprule
Algorithm & Variant & Fault & First & Second & Third \\
\midrule
\endfirsthead
\toprule
Algorithm & Variant & Fault & First & Second & Third \\
\midrule
\endhead
PCMCI & Linear & Fault 1 & \texttt{XMEAS\_1} & \texttt{XMEAS\_20} & \texttt{XMEAS\_4} \\
 & Linear & Fault 2 & \texttt{XMEAS\_10} & \texttt{XMEAS\_30} & \texttt{XMV\_11} \\
 & Linear & Fault 3 & \texttt{XMEAS\_20} & \texttt{XMV\_9} & \texttt{XMV\_5} \\
 & Linear & Fault 4 & \texttt{XMV\_10} & \texttt{XMEAS\_20} & \texttt{XMV\_9} \\
 & Linear & Fault 5 & \texttt{XMV\_11} & \texttt{XMEAS\_17} & \texttt{XMEAS\_20} \\
 & Linear & Fault 6 & \texttt{XMV\_11} & \texttt{XMV\_3} & \texttt{XMV\_10} \\
 & Linear & Fault 7 & \texttt{XMV\_4} & \texttt{XMEAS\_20} & \texttt{XMV\_9} \\
 & Linear & Fault 8 & \texttt{XMEAS\_20} & \texttt{XMV\_5} & \texttt{XMEAS\_16} \\
 & Linear & Fault 9 & \texttt{XMEAS\_20} & \texttt{XMV\_9} & \texttt{XMV\_5} \\
 & Linear & Fault 10 & \texttt{XMEAS\_20} & \texttt{XMV\_9} & \texttt{XMEAS\_27} \\
 & Linear & Fault 11 & \texttt{XMEAS\_20} & \texttt{XMV\_10} & \texttt{XMEAS\_9} \\
 & Linear & Fault 12 & \texttt{XMEAS\_20} & \texttt{XMEAS\_27} & \texttt{XMV\_11} \\
 & Linear & Fault 13 & \texttt{XMEAS\_20} & \texttt{XMV\_9} & \texttt{XMEAS\_19} \\
 & Linear & Fault 14 & \texttt{XMEAS\_9} & \texttt{XMEAS\_21} & \texttt{XMEAS\_20} \\
 & Linear & Fault 15 & \texttt{XMEAS\_20} & \texttt{XMV\_9} & \texttt{XMV\_5} \\
 & Linear & Fault 16 & \texttt{XMV\_9} & \texttt{XMEAS\_20} & \texttt{XMV\_5} \\
 & Linear & Fault 17 & \texttt{XMEAS\_21} & \texttt{XMEAS\_9} & \texttt{XMV\_10} \\
 & Linear & Fault 18 & \texttt{XMV\_5} & \texttt{XMV\_2} & \texttt{XMV\_11} \\
 & Linear & Fault 19 & \texttt{XMEAS\_20} & \texttt{XMEAS\_16} & \texttt{XMV\_9} \\
 & Linear & Fault 20 & \texttt{XMEAS\_13} & \texttt{XMV\_5} & \texttt{XMV\_9} \\
 & Poly-3 & Fault 1 & \texttt{XMEAS\_1} & \texttt{XMEAS\_20} & \texttt{XMEAS\_4} \\
 & Poly-3 & Fault 2 & \texttt{XMEAS\_10} & \texttt{XMEAS\_22} & \texttt{XMEAS\_30} \\
 & Poly-3 & Fault 3 & \texttt{XMEAS\_20} & \texttt{XMEAS\_27} & \texttt{XMV\_9} \\
 & Poly-3 & Fault 4 & \texttt{XMV\_10} & \texttt{XMEAS\_20} & \texttt{XMEAS\_27} \\
 & Poly-3 & Fault 5 & \texttt{XMEAS\_20} & \texttt{XMV\_11} & \texttt{XMEAS\_17} \\
 & Poly-3 & Fault 6 & \texttt{XMV\_11} & \texttt{XMEAS\_7} & \texttt{XMEAS\_20} \\
 & Poly-3 & Fault 7 & \texttt{XMEAS\_20} & \texttt{XMV\_4} & \texttt{XMEAS\_7} \\
 & Poly-3 & Fault 8 & \texttt{XMEAS\_20} & \texttt{XMV\_5} & \texttt{XMEAS\_7} \\
 & Poly-3 & Fault 9 & \texttt{XMEAS\_20} & \texttt{XMEAS\_27} & \texttt{XMV\_9} \\
 & Poly-3 & Fault 10 & \texttt{XMEAS\_20} & \texttt{XMEAS\_18} & \texttt{XMV\_9} \\
 & Poly-3 & Fault 11 & \texttt{XMEAS\_20} & \texttt{XMEAS\_9} & \texttt{XMV\_10} \\
 & Poly-3 & Fault 12 & \texttt{XMEAS\_20} & \texttt{XMEAS\_7} & \texttt{XMEAS\_38} \\
 & Poly-3 & Fault 13 & \texttt{XMEAS\_20} & \texttt{XMEAS\_19} & \texttt{XMEAS\_27} \\
 & Poly-3 & Fault 14 & \texttt{XMEAS\_9} & \texttt{XMEAS\_20} & \texttt{XMEAS\_21} \\
 & Poly-3 & Fault 15 & \texttt{XMEAS\_20} & \texttt{XMEAS\_27} & \texttt{XMEAS\_18} \\
 & Poly-3 & Fault 16 & \texttt{XMV\_9} & \texttt{XMEAS\_20} & \texttt{XMEAS\_19} \\
 & Poly-3 & Fault 17 & \texttt{XMEAS\_21} & \texttt{XMEAS\_9} & \texttt{XMV\_10} \\
 & Poly-3 & Fault 18 & \texttt{XMV\_11} & \texttt{XMEAS\_20} & \texttt{XMV\_2} \\
 & Poly-3 & Fault 19 & \texttt{XMEAS\_20} & \texttt{XMV\_9} & \texttt{XMEAS\_18} \\
 & Poly-3 & Fault 20 & \texttt{XMV\_5} & \texttt{XMEAS\_13} & \texttt{XMEAS\_20} \\
 & RBF & Fault 1 & \texttt{XMEAS\_1} & \texttt{XMEAS\_29} & \texttt{XMEAS\_31} \\
 & RBF & Fault 2 & \texttt{XMEAS\_10} & \texttt{XMEAS\_34} & \texttt{XMEAS\_28} \\
 & RBF & Fault 3 & \texttt{XMEAS\_1} & \texttt{XMEAS\_29} & \texttt{XMEAS\_19} \\
 & RBF & Fault 4 & \texttt{XMEAS\_1} & \texttt{XMV\_10} & \texttt{XMEAS\_29} \\
 & RBF & Fault 5 & \texttt{XMEAS\_1} & \texttt{XMEAS\_19} & \texttt{XMEAS\_18} \\
 & RBF & Fault 6 & \texttt{XMEAS\_25} & \texttt{XMEAS\_29} & \texttt{XMEAS\_7} \\
 & RBF & Fault 7 & \texttt{XMEAS\_1} & \texttt{XMEAS\_19} & \texttt{XMEAS\_7} \\
 & RBF & Fault 8 & \texttt{XMEAS\_1} & \texttt{XMEAS\_29} & \texttt{XMEAS\_10} \\
 & RBF & Fault 9 & \texttt{XMEAS\_1} & \texttt{XMEAS\_29} & \texttt{XMEAS\_19} \\
 & RBF & Fault 10 & \texttt{XMEAS\_1} & \texttt{XMEAS\_18} & \texttt{XMEAS\_11} \\
 & RBF & Fault 11 & \texttt{XMV\_10} & \texttt{XMEAS\_9} & \texttt{XMEAS\_1} \\
 & RBF & Fault 12 & \texttt{XMEAS\_1} & \texttt{XMEAS\_7} & \texttt{XMEAS\_19} \\
 & RBF & Fault 13 & \texttt{XMEAS\_19} & \texttt{XMEAS\_38} & \texttt{XMEAS\_1} \\
 & RBF & Fault 14 & \texttt{XMV\_10} & \texttt{XMEAS\_9} & \texttt{XMEAS\_1} \\
 & RBF & Fault 15 & \texttt{XMEAS\_1} & \texttt{XMEAS\_29} & \texttt{XMEAS\_19} \\
 & RBF & Fault 16 & \texttt{XMEAS\_19} & \texttt{XMEAS\_1} & \texttt{XMEAS\_20} \\
 & RBF & Fault 17 & \texttt{XMEAS\_21} & \texttt{XMEAS\_9} & \texttt{XMV\_10} \\
 & RBF & Fault 18 & \texttt{XMV\_2} & \texttt{XMEAS\_7} & \texttt{XMEAS\_1} \\
 & RBF & Fault 19 & \texttt{XMEAS\_1} & \texttt{XMEAS\_29} & \texttt{XMEAS\_19} \\
 & RBF & Fault 20 & \texttt{XMEAS\_20} & \texttt{XMEAS\_1} & \texttt{XMEAS\_11} \\
GES & Linear & Fault 1 & \texttt{XMV\_3} & \texttt{XMEAS\_4} & \texttt{XMV\_5} \\
 & Linear & Fault 2 & \texttt{XMEAS\_4} & \texttt{XMEAS\_16} & \texttt{XMV\_6} \\
 & Linear & Fault 3 & \texttt{XMV\_5} & \texttt{XMEAS\_21} & \texttt{XMEAS\_18} \\
 & Linear & Fault 4 & \texttt{XMV\_10} & \texttt{XMV\_5} & \texttt{XMEAS\_18} \\
 & Linear & Fault 5 & \texttt{XMV\_5} & \texttt{XMV\_9} & \texttt{XMEAS\_21} \\
 & Linear & Fault 6 & \texttt{XMV\_10} & \texttt{XMEAS\_20} & \texttt{XMV\_3} \\
 & Linear & Fault 7 & \texttt{XMV\_5} & \texttt{XMV\_9} & \texttt{XMEAS\_21} \\
 & Linear & Fault 8 & \texttt{XMV\_5} & \texttt{XMEAS\_34} & \texttt{XMEAS\_31} \\
 & Linear & Fault 9 & \texttt{XMV\_5} & \texttt{XMEAS\_18} & \texttt{XMV\_6} \\
 & Linear & Fault 10 & \texttt{XMEAS\_18} & \texttt{XMV\_5} & \texttt{XMEAS\_20} \\
 & Linear & Fault 11 & \texttt{XMV\_10} & \texttt{XMV\_5} & \texttt{XMV\_6} \\
 & Linear & Fault 12 & \texttt{XMV\_5} & \texttt{XMV\_9} & \texttt{XMEAS\_21} \\
 & Linear & Fault 13 & \texttt{XMV\_5} & \texttt{XMEAS\_18} & \texttt{XMV\_9} \\
 & Linear & Fault 14 & \texttt{XMEAS\_21} & \texttt{XMV\_10} & \texttt{XMV\_5} \\
 & Linear & Fault 15 & \texttt{XMV\_5} & \texttt{XMEAS\_18} & \texttt{XMV\_6} \\
 & Linear & Fault 16 & \texttt{XMV\_9} & \texttt{XMV\_5} & \texttt{XMEAS\_18} \\
 & Linear & Fault 17 & \texttt{XMEAS\_21} & \texttt{XMV\_10} & \texttt{XMV\_5} \\
 & Linear & Fault 18 & \texttt{XMEAS\_21} & \texttt{XMV\_5} & \texttt{XMV\_1} \\
 & Linear & Fault 19 & \texttt{XMV\_5} & \texttt{XMEAS\_20} & \texttt{XMEAS\_18} \\
 & Linear & Fault 20 & \texttt{XMV\_5} & \texttt{XMEAS\_13} & \texttt{XMEAS\_20} \\
 & Poly-3 & Fault 1 & \texttt{XMV\_3} & \texttt{XMEAS\_4} & \texttt{XMEAS\_16} \\
 & Poly-3 & Fault 2 & \texttt{XMEAS\_4} & \texttt{XMEAS\_16} & \texttt{XMV\_3} \\
 & Poly-3 & Fault 3 & \texttt{XMV\_9} & \texttt{XMEAS\_21} & \texttt{XMEAS\_13} \\
 & Poly-3 & Fault 4 & \texttt{XMV\_10} & \texttt{XMV\_9} & \texttt{XMEAS\_13} \\
 & Poly-3 & Fault 5 & \texttt{XMV\_5} & \texttt{XMEAS\_21} & \texttt{XMV\_9} \\
 & Poly-3 & Fault 6 & \texttt{XMV\_10} & \texttt{XMV\_3} & \texttt{XMEAS\_20} \\
 & Poly-3 & Fault 7 & \texttt{XMV\_5} & \texttt{XMEAS\_33} & \texttt{XMEAS\_21} \\
 & Poly-3 & Fault 8 & \texttt{XMEAS\_34} & \texttt{XMEAS\_31} & \texttt{XMV\_9} \\
 & Poly-3 & Fault 9 & \texttt{XMV\_9} & \texttt{XMEAS\_13} & \texttt{XMEAS\_21} \\
 & Poly-3 & Fault 10 & \texttt{XMV\_9} & \texttt{XMEAS\_18} & \texttt{XMEAS\_7} \\
 & Poly-3 & Fault 11 & \texttt{XMV\_10} & \texttt{XMV\_9} & \texttt{XMEAS\_13} \\
 & Poly-3 & Fault 12 & \texttt{XMV\_5} & \texttt{XMEAS\_7} & \texttt{XMV\_9} \\
 & Poly-3 & Fault 13 & \texttt{XMEAS\_38} & \texttt{XMV\_9} & \texttt{XMV\_5} \\
 & Poly-3 & Fault 14 & \texttt{XMEAS\_21} & \texttt{XMV\_10} & \texttt{XMV\_9} \\
 & Poly-3 & Fault 15 & \texttt{XMV\_9} & \texttt{XMEAS\_21} & \texttt{XMEAS\_13} \\
 & Poly-3 & Fault 16 & \texttt{XMV\_9} & \texttt{XMEAS\_18} & \texttt{XMEAS\_33} \\
 & Poly-3 & Fault 17 & \texttt{XMEAS\_21} & \texttt{XMV\_10} & \texttt{XMV\_9} \\
 & Poly-3 & Fault 18 & \texttt{XMEAS\_21} & \texttt{XMEAS\_7} & \texttt{XMV\_5} \\
 & Poly-3 & Fault 19 & \texttt{XMV\_9} & \texttt{XMEAS\_21} & \texttt{XMEAS\_17} \\
 & Poly-3 & Fault 20 & \texttt{XMEAS\_13} & \texttt{XMV\_5} & \texttt{XMEAS\_20} \\
 & RBF & Fault 1 & \texttt{XMEAS\_1} & \texttt{XMV\_3} & \texttt{XMEAS\_4} \\
 & RBF & Fault 2 & \texttt{XMV\_6} & \texttt{XMEAS\_34} & \texttt{XMV\_9} \\
 & RBF & Fault 3 & \texttt{XMEAS\_21} & \texttt{XMEAS\_38} & \texttt{XMV\_9} \\
 & RBF & Fault 4 & \texttt{XMV\_10} & \texttt{XMEAS\_38} & \texttt{XMV\_9} \\
 & RBF & Fault 5 & \texttt{XMEAS\_38} & \texttt{XMV\_9} & \texttt{XMEAS\_18} \\
 & RBF & Fault 6 & \texttt{XMV\_10} & \texttt{XMEAS\_16} & \texttt{XMV\_5} \\
 & RBF & Fault 7 & \texttt{XMEAS\_38} & \texttt{XMV\_9} & \texttt{XMEAS\_18} \\
 & RBF & Fault 8 & \texttt{XMEAS\_1} & \texttt{XMEAS\_38} & \texttt{XMV\_6} \\
 & RBF & Fault 9 & \texttt{XMEAS\_38} & \texttt{XMV\_9} & \texttt{XMEAS\_18} \\
 & RBF & Fault 10 & \texttt{XMEAS\_18} & \texttt{XMEAS\_38} & \texttt{XMV\_9} \\
 & RBF & Fault 11 & \texttt{XMV\_10} & \texttt{XMEAS\_21} & \texttt{XMEAS\_38} \\
 & RBF & Fault 12 & \texttt{XMEAS\_38} & \texttt{XMEAS\_20} & \texttt{XMV\_5} \\
 & RBF & Fault 13 & \texttt{XMV\_9} & \texttt{XMEAS\_18} & \texttt{XMEAS\_13} \\
 & RBF & Fault 14 & \texttt{XMV\_10} & \texttt{XMEAS\_21} & \texttt{XMEAS\_38} \\
 & RBF & Fault 15 & \texttt{XMEAS\_38} & \texttt{XMV\_9} & \texttt{XMEAS\_18} \\
 & RBF & Fault 16 & \texttt{XMEAS\_38} & \texttt{XMEAS\_18} & \texttt{XMV\_9} \\
 & RBF & Fault 17 & \texttt{XMEAS\_21} & \texttt{XMV\_10} & \texttt{XMV\_1} \\
 & RBF & Fault 18 & \texttt{XMV\_1} & \texttt{XMEAS\_21} & \texttt{XMEAS\_4} \\
 & RBF & Fault 19 & \texttt{XMEAS\_38} & \texttt{XMV\_5} & \texttt{XMV\_9} \\
 & RBF & Fault 20 & \texttt{XMV\_5} & \texttt{XMEAS\_20} & \texttt{XMEAS\_13} \\
CAM-UV & Linear & Fault 1 & \texttt{XMEAS\_1} & \texttt{XMEAS\_19} & \texttt{XMEAS\_20} \\
 & Linear & Fault 2 & \texttt{XMEAS\_28} & \texttt{XMEAS\_10} & \texttt{XMEAS\_19} \\
 & Linear & Fault 3 & \texttt{XMEAS\_20} & \texttt{XMEAS\_21} & \texttt{XMV\_6} \\
 & Linear & Fault 4 & \texttt{XMV\_10} & \texttt{XMEAS\_9} & \texttt{XMEAS\_20} \\
 & Linear & Fault 5 & \texttt{XMEAS\_20} & \texttt{XMV\_11} & \texttt{XMEAS\_19} \\
 & Linear & Fault 6 & \texttt{XMEAS\_21} & \texttt{XMEAS\_28} & \texttt{XMV\_5} \\
 & Linear & Fault 7 & \texttt{XMV\_4} & \texttt{XMEAS\_20} & \texttt{XMEAS\_19} \\
 & Linear & Fault 8 & \texttt{XMEAS\_28} & \texttt{XMEAS\_20} & \texttt{XMEAS\_10} \\
 & Linear & Fault 9 & \texttt{XMEAS\_20} & \texttt{XMV\_6} & \texttt{XMEAS\_28} \\
 & Linear & Fault 10 & \texttt{XMEAS\_20} & \texttt{XMEAS\_18} & \texttt{XMEAS\_38} \\
 & Linear & Fault 11 & \texttt{XMV\_10} & \texttt{XMEAS\_20} & \texttt{XMV\_6} \\
 & Linear & Fault 12 & \texttt{XMEAS\_20} & \texttt{XMEAS\_16} & \texttt{XMEAS\_19} \\
 & Linear & Fault 13 & \texttt{XMEAS\_19} & \texttt{XMEAS\_20} & \texttt{XMEAS\_16} \\
 & Linear & Fault 14 & \texttt{XMEAS\_21} & \texttt{XMEAS\_9} & \texttt{XMV\_10} \\
 & Linear & Fault 15 & \texttt{XMEAS\_20} & \texttt{XMV\_6} & \texttt{XMEAS\_28} \\
 & Linear & Fault 16 & \texttt{XMEAS\_20} & \texttt{XMV\_6} & \texttt{XMEAS\_19} \\
 & Linear & Fault 17 & \texttt{XMEAS\_21} & \texttt{XMV\_10} & \texttt{XMEAS\_9} \\
 & Linear & Fault 18 & \texttt{XMV\_1} & \texttt{XMV\_5} & \texttt{XMV\_10} \\
 & Linear & Fault 19 & \texttt{XMV\_6} & \texttt{XMV\_5} & \texttt{XMEAS\_28} \\
 & Linear & Fault 20 & \texttt{XMV\_5} & \texttt{XMEAS\_20} & \texttt{XMEAS\_13} \\
 & Poly-3 & Fault 1 & \texttt{XMEAS\_1} & \texttt{XMEAS\_20} & \texttt{XMV\_3} \\
 & Poly-3 & Fault 2 & \texttt{XMEAS\_28} & \texttt{XMEAS\_10} & \texttt{XMEAS\_19} \\
 & Poly-3 & Fault 3 & \texttt{XMEAS\_38} & \texttt{XMEAS\_20} & \texttt{XMEAS\_1} \\
 & Poly-3 & Fault 4 & \texttt{XMV\_10} & \texttt{XMEAS\_9} & \texttt{XMEAS\_20} \\
 & Poly-3 & Fault 5 & \texttt{XMEAS\_20} & \texttt{XMV\_11} & \texttt{XMEAS\_18} \\
 & Poly-3 & Fault 6 & \texttt{XMV\_10} & \texttt{XMEAS\_21} & \texttt{XMEAS\_23} \\
 & Poly-3 & Fault 7 & \texttt{XMV\_4} & \texttt{XMEAS\_20} & \texttt{XMEAS\_16} \\
 & Poly-3 & Fault 8 & \texttt{XMEAS\_29} & \texttt{XMEAS\_20} & \texttt{XMEAS\_34} \\
 & Poly-3 & Fault 9 & \texttt{XMEAS\_38} & \texttt{XMEAS\_20} & \texttt{XMEAS\_1} \\
 & Poly-3 & Fault 10 & \texttt{XMEAS\_18} & \texttt{XMEAS\_20} & \texttt{XMEAS\_38} \\
 & Poly-3 & Fault 11 & \texttt{XMV\_10} & \texttt{XMEAS\_38} & \texttt{XMEAS\_20} \\
 & Poly-3 & Fault 12 & \texttt{XMEAS\_20} & \texttt{XMEAS\_18} & \texttt{XMEAS\_21} \\
 & Poly-3 & Fault 13 & \texttt{XMEAS\_38} & \texttt{XMEAS\_16} & \texttt{XMEAS\_20} \\
 & Poly-3 & Fault 14 & \texttt{XMEAS\_21} & \texttt{XMV\_10} & \texttt{XMEAS\_9} \\
 & Poly-3 & Fault 15 & \texttt{XMEAS\_20} & \texttt{XMEAS\_38} & \texttt{XMV\_6} \\
 & Poly-3 & Fault 16 & \texttt{XMEAS\_20} & \texttt{XMEAS\_38} & \texttt{XMV\_9} \\
 & Poly-3 & Fault 17 & \texttt{XMEAS\_21} & \texttt{XMEAS\_9} & \texttt{XMEAS\_38} \\
 & Poly-3 & Fault 18 & \texttt{XMEAS\_20} & \texttt{XMV\_1} & \texttt{XMV\_5} \\
 & Poly-3 & Fault 19 & \texttt{XMV\_5} & \texttt{XMEAS\_38} & \texttt{XMEAS\_20} \\
 & Poly-3 & Fault 20 & \texttt{XMV\_5} & \texttt{XMEAS\_20} & \texttt{XMEAS\_13} \\
 & RBF & Fault 1 & \texttt{XMEAS\_1} & \texttt{XMV\_3} & \texttt{XMV\_9} \\
 & RBF & Fault 2 & \texttt{XMEAS\_10} & \texttt{XMV\_6} & \texttt{XMEAS\_28} \\
 & RBF & Fault 3 & \texttt{XMEAS\_21} & \texttt{XMEAS\_14} & \texttt{XMV\_9} \\
 & RBF & Fault 4 & \texttt{XMV\_10} & \texttt{XMEAS\_9} & \texttt{XMEAS\_14} \\
 & RBF & Fault 5 & \texttt{XMEAS\_18} & \texttt{XMV\_9} & \texttt{XMEAS\_38} \\
 & RBF & Fault 6 & \texttt{XMV\_10} & \texttt{XMEAS\_16} & \texttt{XMV\_3} \\
 & RBF & Fault 7 & \texttt{XMEAS\_38} & \texttt{XMV\_9} & \texttt{XMEAS\_18} \\
 & RBF & Fault 8 & \texttt{XMEAS\_38} & \texttt{XMV\_3} & \texttt{XMEAS\_10} \\
 & RBF & Fault 9 & \texttt{XMEAS\_38} & \texttt{XMEAS\_14} & \texttt{XMV\_9} \\
 & RBF & Fault 10 & \texttt{XMEAS\_18} & \texttt{XMV\_9} & \texttt{XMEAS\_38} \\
 & RBF & Fault 11 & \texttt{XMEAS\_9} & \texttt{XMV\_10} & \texttt{XMEAS\_38} \\
 & RBF & Fault 12 & \texttt{XMEAS\_38} & \texttt{XMEAS\_13} & \texttt{XMEAS\_7} \\
 & RBF & Fault 13 & \texttt{XMV\_9} & \texttt{XMEAS\_18} & \texttt{XMEAS\_13} \\
 & RBF & Fault 14 & \texttt{XMEAS\_9} & \texttt{XMV\_10} & \texttt{XMEAS\_21} \\
 & RBF & Fault 15 & \texttt{XMEAS\_38} & \texttt{XMV\_9} & \texttt{XMEAS\_18} \\
 & RBF & Fault 16 & \texttt{XMEAS\_38} & \texttt{XMEAS\_18} & \texttt{XMV\_9} \\
 & RBF & Fault 17 & \texttt{XMEAS\_21} & \texttt{XMV\_10} & \texttt{XMEAS\_9} \\
 & RBF & Fault 18 & \texttt{XMEAS\_20} & \texttt{XMEAS\_21} & \texttt{XMEAS\_13} \\
 & RBF & Fault 19 & \texttt{XMEAS\_38} & \texttt{XMEAS\_18} & \texttt{XMV\_9} \\
 & RBF & Fault 20 & \texttt{XMV\_5} & \texttt{XMEAS\_20} & \texttt{XMEAS\_13} \\
\bottomrule
\end{longtable}}


\section{Per-Configuration Attribution Plots}

Figures~\ref{fig:pcmci-robotic-linear-ridge} to~\ref{fig:camuv-industrial-rbf-ridge} give one figure per configuration, showing the ranked attribution for each attack class in that domain. These are the graphical form of the tables above, included so that the relative magnitudes and not only the ordering are visible.

\begin{figure}[htbp]
  \centering
  \includegraphics[width=0.85\textwidth,height=0.78\textheight,keepaspectratio]{pcmci_robotic_linear_ridge.png}
  \caption[Attribution: PCMCI, Robotics, linear/ridge]{Top attributed variables per attack class. PCMCI, Robotics domain, linear/ridge.}
  \label{fig:pcmci-robotic-linear-ridge}
\end{figure}

\begin{figure}[htbp]
  \centering
  \includegraphics[width=0.85\textwidth,height=0.78\textheight,keepaspectratio]{pcmci_robotic_polynomial_degree_3.png}
  \caption[Attribution: PCMCI, Robotics, polynomial degree 3]{Top attributed variables per attack class. PCMCI, Robotics domain, polynomial degree 3.}
  \label{fig:pcmci-robotic-polynomial-degree-3}
\end{figure}

\begin{figure}[htbp]
  \centering
  \includegraphics[width=0.85\textwidth,height=0.78\textheight,keepaspectratio]{pcmci_robotic_rbf_ridge.png}
  \caption[Attribution: PCMCI, Robotics, RBF + ridge]{Top attributed variables per attack class. PCMCI, Robotics domain, RBF + ridge.}
  \label{fig:pcmci-robotic-rbf-ridge}
\end{figure}

\begin{figure}[htbp]
  \centering
  \includegraphics[width=0.85\textwidth,height=0.78\textheight,keepaspectratio]{pcmci_healthcare_linear_ridge.png}
  \caption[Attribution: PCMCI, Healthcare, linear/ridge]{Top attributed variables per attack class. PCMCI, Healthcare domain, linear/ridge.}
  \label{fig:pcmci-healthcare-linear-ridge}
\end{figure}

\begin{figure}[htbp]
  \centering
  \includegraphics[width=0.85\textwidth,height=0.78\textheight,keepaspectratio]{pcmci_healthcare_polynomial_degree_3.png}
  \caption[Attribution: PCMCI, Healthcare, polynomial degree 3]{Top attributed variables per attack class. PCMCI, Healthcare domain, polynomial degree 3.}
  \label{fig:pcmci-healthcare-polynomial-degree-3}
\end{figure}

\begin{figure}[htbp]
  \centering
  \includegraphics[width=0.85\textwidth,height=0.78\textheight,keepaspectratio]{pcmci_healthcare_rbf_ridge.png}
  \caption[Attribution: PCMCI, Healthcare, RBF + ridge]{Top attributed variables per attack class. PCMCI, Healthcare domain, RBF + ridge.}
  \label{fig:pcmci-healthcare-rbf-ridge}
\end{figure}

\begin{figure}[htbp]
  \centering
  \includegraphics[width=0.85\textwidth,height=0.78\textheight,keepaspectratio]{pcmci_industrial_linear_ridge.png}
  \caption[Attribution: PCMCI, Industrial, linear/ridge]{Top attributed variables per attack class. PCMCI, Industrial domain, linear/ridge.}
  \label{fig:pcmci-industrial-linear-ridge}
\end{figure}

\begin{figure}[htbp]
  \centering
  \includegraphics[width=0.85\textwidth,height=0.78\textheight,keepaspectratio]{pcmci_industrial_polynomial_degree_3.png}
  \caption[Attribution: PCMCI, Industrial, polynomial degree 3]{Top attributed variables per attack class. PCMCI, Industrial domain, polynomial degree 3.}
  \label{fig:pcmci-industrial-polynomial-degree-3}
\end{figure}

\begin{figure}[htbp]
  \centering
  \includegraphics[width=0.85\textwidth,height=0.78\textheight,keepaspectratio]{pcmci_industrial_rbf_ridge.png}
  \caption[Attribution: PCMCI, Industrial, RBF + ridge]{Top attributed variables per attack class. PCMCI, Industrial domain, RBF + ridge.}
  \label{fig:pcmci-industrial-rbf-ridge}
\end{figure}

\begin{figure}[htbp]
  \centering
  \includegraphics[width=0.85\textwidth,height=0.78\textheight,keepaspectratio]{ges_robotic_linear_ridge.png}
  \caption[Attribution: GES, Robotics, linear/ridge]{Top attributed variables per attack class. GES, Robotics domain, linear/ridge.}
  \label{fig:ges-robotic-linear-ridge}
\end{figure}

\begin{figure}[htbp]
  \centering
  \includegraphics[width=0.85\textwidth,height=0.78\textheight,keepaspectratio]{ges_robotic_polynomial_degree_3.png}
  \caption[Attribution: GES, Robotics, polynomial degree 3]{Top attributed variables per attack class. GES, Robotics domain, polynomial degree 3.}
  \label{fig:ges-robotic-polynomial-degree-3}
\end{figure}

\begin{figure}[htbp]
  \centering
  \includegraphics[width=0.85\textwidth,height=0.78\textheight,keepaspectratio]{ges_robotic_rbf_ridge.png}
  \caption[Attribution: GES, Robotics, RBF + ridge]{Top attributed variables per attack class. GES, Robotics domain, RBF + ridge.}
  \label{fig:ges-robotic-rbf-ridge}
\end{figure}

\begin{figure}[htbp]
  \centering
  \includegraphics[width=0.85\textwidth,height=0.78\textheight,keepaspectratio]{ges_healthcare_linear_ridge.png}
  \caption[Attribution: GES, Healthcare, linear/ridge]{Top attributed variables per attack class. GES, Healthcare domain, linear/ridge.}
  \label{fig:ges-healthcare-linear-ridge}
\end{figure}

\begin{figure}[htbp]
  \centering
  \includegraphics[width=0.85\textwidth,height=0.78\textheight,keepaspectratio]{ges_healthcare_polynomial_degree_3.png}
  \caption[Attribution: GES, Healthcare, polynomial degree 3]{Top attributed variables per attack class. GES, Healthcare domain, polynomial degree 3.}
  \label{fig:ges-healthcare-polynomial-degree-3}
\end{figure}

\begin{figure}[htbp]
  \centering
  \includegraphics[width=0.85\textwidth,height=0.78\textheight,keepaspectratio]{ges_healthcare_rbf_ridge.png}
  \caption[Attribution: GES, Healthcare, RBF + ridge]{Top attributed variables per attack class. GES, Healthcare domain, RBF + ridge.}
  \label{fig:ges-healthcare-rbf-ridge}
\end{figure}

\begin{figure}[htbp]
  \centering
  \includegraphics[width=0.85\textwidth,height=0.78\textheight,keepaspectratio]{ges_industrial_linear_ridge.png}
  \caption[Attribution: GES, Industrial, linear/ridge]{Top attributed variables per attack class. GES, Industrial domain, linear/ridge.}
  \label{fig:ges-industrial-linear-ridge}
\end{figure}

\begin{figure}[htbp]
  \centering
  \includegraphics[width=0.85\textwidth,height=0.78\textheight,keepaspectratio]{ges_industrial_polynomial_degree_3.png}
  \caption[Attribution: GES, Industrial, polynomial degree 3]{Top attributed variables per attack class. GES, Industrial domain, polynomial degree 3.}
  \label{fig:ges-industrial-polynomial-degree-3}
\end{figure}

\begin{figure}[htbp]
  \centering
  \includegraphics[width=0.85\textwidth,height=0.78\textheight,keepaspectratio]{ges_industrial_rbf_ridge.png}
  \caption[Attribution: GES, Industrial, RBF + ridge]{Top attributed variables per attack class. GES, Industrial domain, RBF + ridge.}
  \label{fig:ges-industrial-rbf-ridge}
\end{figure}

\begin{figure}[htbp]
  \centering
  \includegraphics[width=0.85\textwidth,height=0.78\textheight,keepaspectratio]{camuv_robotics_linear_ridge.png}
  \caption[Attribution: CAM-UV, Robotics, linear/ridge]{Top attributed variables per attack class. CAM-UV, Robotics domain, linear/ridge.}
  \label{fig:camuv-robotics-linear-ridge}
\end{figure}

\begin{figure}[htbp]
  \centering
  \includegraphics[width=0.85\textwidth,height=0.78\textheight,keepaspectratio]{camuv_robotics_polynomial_degree_3.png}
  \caption[Attribution: CAM-UV, Robotics, polynomial degree 3]{Top attributed variables per attack class. CAM-UV, Robotics domain, polynomial degree 3.}
  \label{fig:camuv-robotics-polynomial-degree-3}
\end{figure}

\begin{figure}[htbp]
  \centering
  \includegraphics[width=0.85\textwidth,height=0.78\textheight,keepaspectratio]{camuv_robotics_rbf_ridge.png}
  \caption[Attribution: CAM-UV, Robotics, RBF + ridge]{Top attributed variables per attack class. CAM-UV, Robotics domain, RBF + ridge.}
  \label{fig:camuv-robotics-rbf-ridge}
\end{figure}

\begin{figure}[htbp]
  \centering
  \includegraphics[width=0.85\textwidth,height=0.78\textheight,keepaspectratio]{camuv_healthcare_linear_ridge.png}
  \caption[Attribution: CAM-UV, Healthcare, linear/ridge]{Top attributed variables per attack class. CAM-UV, Healthcare domain, linear/ridge.}
  \label{fig:camuv-healthcare-linear-ridge}
\end{figure}

\begin{figure}[htbp]
  \centering
  \includegraphics[width=0.85\textwidth,height=0.78\textheight,keepaspectratio]{camuv_healthcare_polynomial_degree_3.png}
  \caption[Attribution: CAM-UV, Healthcare, polynomial degree 3]{Top attributed variables per attack class. CAM-UV, Healthcare domain, polynomial degree 3.}
  \label{fig:camuv-healthcare-polynomial-degree-3}
\end{figure}

\begin{figure}[htbp]
  \centering
  \includegraphics[width=0.85\textwidth,height=0.78\textheight,keepaspectratio]{camuv_healthcare_rbf_ridge.png}
  \caption[Attribution: CAM-UV, Healthcare, RBF + ridge]{Top attributed variables per attack class. CAM-UV, Healthcare domain, RBF + ridge.}
  \label{fig:camuv-healthcare-rbf-ridge}
\end{figure}

\begin{figure}[htbp]
  \centering
  \includegraphics[width=0.85\textwidth,height=0.78\textheight,keepaspectratio]{camuv_industrial_linear_ridge.png}
  \caption[Attribution: CAM-UV, Industrial, linear/ridge]{Top attributed variables per attack class. CAM-UV, Industrial domain, linear/ridge.}
  \label{fig:camuv-industrial-linear-ridge}
\end{figure}

\begin{figure}[htbp]
  \centering
  \includegraphics[width=0.85\textwidth,height=0.78\textheight,keepaspectratio]{camuv_industrial_polynomial_degree_3.png}
  \caption[Attribution: CAM-UV, Industrial, polynomial degree 3]{Top attributed variables per attack class. CAM-UV, Industrial domain, polynomial degree 3.}
  \label{fig:camuv-industrial-polynomial-degree-3}
\end{figure}

\begin{figure}[htbp]
  \centering
  \includegraphics[width=0.85\textwidth,height=0.78\textheight,keepaspectratio]{camuv_industrial_rbf_ridge.png}
  \caption[Attribution: CAM-UV, Industrial, RBF + ridge]{Top attributed variables per attack class. CAM-UV, Industrial domain, RBF + ridge.}
  \label{fig:camuv-industrial-rbf-ridge}
\end{figure}

% Flush the per-configuration figures before the next section starts. They are
% full-page floats, so without this LaTeX defers them past two section breaks and
% prints them inside the baseline-attribution section they do not belong to.
\clearpage

\section{Paired AUC Comparisons}

Every comparison DeLong's test can reach, all 54 of them. The 54 it cannot reach are the pairs in which PCMCI faces one of the other two algorithms, and they are absent for the reason set out in Section~\ref{sec:uncertainty} rather than because they were not attempted.

% GENERATED by Thesis/tools/make_tables.py -- do not edit by hand.
% Every figure below is read from a CSV written by the pipelines.
{\footnotesize
\begin{longtable}{llrrrrr}
\caption[Every testable paired comparison, by DeLong's test]{Every testable paired comparison, by DeLong's test. $\Delta$ is AUC$_A$ minus AUC$_B$; $p$ is two-sided and is set in bold where it clears $\alpha = 0.05$. Pairs absent from this table are the 54 that share no sample set; the tool that writes this table also writes those, with the reason each one fails, to a companion CSV.}\label{tab:delong-full}\\
\toprule
A & B & AUC$_A$ & AUC$_B$ & $\Delta$ & $z$ & $p$ \\
\midrule
\endfirsthead
\toprule
A & B & AUC$_A$ & AUC$_B$ & $\Delta$ & $z$ & $p$ \\
\midrule
\endhead
\multicolumn{7}{l}{\textit{Healthcare (ECG)}} \\
CAM-UV Linear & CAM-UV Poly-3 & 0.7861 & 0.7691 & 0.0170 & 2.1897 & \textbf{0.028544} \\
CAM-UV Linear & CAM-UV RBF & 0.7861 & 0.7692 & 0.0169 & 2.5283 & \textbf{0.011461} \\
CAM-UV Linear & GES Linear & 0.7861 & 0.8114 & -0.0254 & -5.5553 & \textbf{0.000000} \\
CAM-UV Linear & GES Poly-3 & 0.7861 & 0.7867 & -0.0006 & -0.0853 & 0.932044 \\
CAM-UV Linear & GES RBF & 0.7861 & 0.7620 & 0.0241 & 3.4719 & \textbf{0.000517} \\
CAM-UV Poly-3 & CAM-UV RBF & 0.7691 & 0.7692 & -0.0001 & -0.0143 & 0.988601 \\
CAM-UV Poly-3 & GES Linear & 0.7691 & 0.8114 & -0.0423 & -5.2519 & \textbf{0.000000} \\
CAM-UV Poly-3 & GES Poly-3 & 0.7691 & 0.7867 & -0.0176 & -3.9817 & \textbf{0.000068} \\
CAM-UV Poly-3 & GES RBF & 0.7691 & 0.7620 & 0.0071 & 0.9705 & 0.331804 \\
CAM-UV RBF & GES Linear & 0.7692 & 0.8114 & -0.0422 & -6.1599 & \textbf{0.000000} \\
CAM-UV RBF & GES Poly-3 & 0.7692 & 0.7867 & -0.0175 & -2.6142 & \textbf{0.008944} \\
CAM-UV RBF & GES RBF & 0.7692 & 0.7620 & 0.0072 & 1.9235 & 0.054424 \\
\midrule
\multicolumn{7}{l}{\textit{Industrial (TEP)}} \\
CAM-UV Linear & CAM-UV Poly-3 & 0.8428 & 0.8212 & 0.0217 & 1.8245 & 0.068083 \\
CAM-UV Linear & CAM-UV RBF & 0.8428 & 0.8922 & -0.0494 & -6.3027 & \textbf{0.000000} \\
CAM-UV Linear & GES Linear & 0.8428 & 0.8407 & 0.0021 & 0.2920 & 0.770303 \\
CAM-UV Linear & GES Poly-3 & 0.8428 & 0.8849 & -0.0421 & -2.5705 & \textbf{0.010154} \\
CAM-UV Linear & GES RBF & 0.8428 & 0.8994 & -0.0565 & -7.1960 & \textbf{0.000000} \\
CAM-UV Poly-3 & CAM-UV RBF & 0.8212 & 0.8922 & -0.0711 & -4.8854 & \textbf{0.000001} \\
CAM-UV Poly-3 & GES Linear & 0.8212 & 0.8407 & -0.0195 & -1.3578 & 0.174527 \\
CAM-UV Poly-3 & GES Poly-3 & 0.8212 & 0.8849 & -0.0638 & -3.0470 & \textbf{0.002311} \\
CAM-UV Poly-3 & GES RBF & 0.8212 & 0.8994 & -0.0782 & -5.4280 & \textbf{0.000000} \\
CAM-UV RBF & GES Linear & 0.8922 & 0.8407 & 0.0515 & 8.0376 & \textbf{0.000000} \\
CAM-UV RBF & GES Poly-3 & 0.8922 & 0.8849 & 0.0073 & 0.5288 & 0.596928 \\
CAM-UV RBF & GES RBF & 0.8922 & 0.8994 & -0.0071 & -1.8047 & 0.071127 \\
\midrule
\multicolumn{7}{l}{\textit{Robotics (Pepper)}} \\
CAM-UV Linear & CAM-UV Poly-3 & 0.6497 & 0.6888 & -0.0391 & -2.4488 & \textbf{0.014334} \\
CAM-UV Linear & CAM-UV RBF & 0.6497 & 0.7818 & -0.1321 & -13.5592 & \textbf{0.000000} \\
CAM-UV Linear & GES Linear & 0.6497 & 0.5906 & 0.0591 & 4.8003 & \textbf{0.000002} \\
CAM-UV Linear & GES Poly-3 & 0.6497 & 0.6999 & -0.0502 & -3.0099 & \textbf{0.002614} \\
CAM-UV Linear & GES RBF & 0.6497 & 0.6516 & -0.0019 & -0.2423 & 0.808520 \\
CAM-UV Poly-3 & CAM-UV RBF & 0.6888 & 0.7818 & -0.0930 & -5.8060 & \textbf{0.000000} \\
CAM-UV Poly-3 & GES Linear & 0.6888 & 0.5906 & 0.0982 & 5.9541 & \textbf{0.000000} \\
CAM-UV Poly-3 & GES Poly-3 & 0.6888 & 0.6999 & -0.0111 & -1.0631 & 0.287753 \\
CAM-UV Poly-3 & GES RBF & 0.6888 & 0.6516 & 0.0372 & 2.5289 & \textbf{0.011443} \\
CAM-UV RBF & GES Linear & 0.7818 & 0.5906 & 0.1912 & 14.4240 & \textbf{0.000000} \\
CAM-UV RBF & GES Poly-3 & 0.7818 & 0.6999 & 0.0819 & 4.9478 & \textbf{0.000001} \\
CAM-UV RBF & GES RBF & 0.7818 & 0.6516 & 0.1302 & 15.3433 & \textbf{0.000000} \\
\midrule
\multicolumn{7}{l}{\textit{Healthcare (ECG)}} \\
GES Linear & GES Poly-3 & 0.8114 & 0.7867 & 0.0247 & 3.2587 & \textbf{0.001119} \\
GES Linear & GES RBF & 0.8114 & 0.7620 & 0.0494 & 7.4872 & \textbf{0.000000} \\
GES Poly-3 & GES RBF & 0.7867 & 0.7620 & 0.0247 & 3.6488 & \textbf{0.000263} \\
\midrule
\multicolumn{7}{l}{\textit{Industrial (TEP)}} \\
GES Linear & GES Poly-3 & 0.8407 & 0.8849 & -0.0442 & -2.7951 & \textbf{0.005188} \\
GES Linear & GES RBF & 0.8407 & 0.8994 & -0.0587 & -8.7444 & \textbf{0.000000} \\
GES Poly-3 & GES RBF & 0.8849 & 0.8994 & -0.0145 & -1.0202 & 0.307616 \\
\midrule
\multicolumn{7}{l}{\textit{Robotics (Pepper)}} \\
GES Linear & GES Poly-3 & 0.5906 & 0.6999 & -0.1093 & -8.2800 & \textbf{0.000000} \\
GES Linear & GES RBF & 0.5906 & 0.6516 & -0.0610 & -4.7679 & \textbf{0.000002} \\
GES Poly-3 & GES RBF & 0.6999 & 0.6516 & 0.0483 & 2.9049 & \textbf{0.003674} \\
\midrule
\multicolumn{7}{l}{\textit{Healthcare (ECG)}} \\
PCMCI Linear & PCMCI Poly-3 & 0.6181 & 0.5075 & 0.1106 & 13.2173 & \textbf{0.000000} \\
PCMCI Linear & PCMCI RBF & 0.6181 & 0.7147 & -0.0966 & -13.0939 & \textbf{0.000000} \\
PCMCI Poly-3 & PCMCI RBF & 0.5075 & 0.7147 & -0.2071 & -22.9203 & \textbf{0.000000} \\
\midrule
\multicolumn{7}{l}{\textit{Industrial (TEP)}} \\
PCMCI Linear & PCMCI Poly-3 & 0.8860 & 0.9151 & -0.0291 & -9.7046 & \textbf{0.000000} \\
PCMCI Linear & PCMCI RBF & 0.8860 & 0.8495 & 0.0365 & 8.1254 & \textbf{0.000000} \\
PCMCI Poly-3 & PCMCI RBF & 0.9151 & 0.8495 & 0.0656 & 17.5271 & \textbf{0.000000} \\
\midrule
\multicolumn{7}{l}{\textit{Robotics (Pepper)}} \\
PCMCI Linear & PCMCI Poly-3 & 0.7825 & 0.8053 & -0.0228 & -0.8914 & 0.372739 \\
PCMCI Linear & PCMCI RBF & 0.7825 & 0.8752 & -0.0927 & -4.3958 & \textbf{0.000011} \\
PCMCI Poly-3 & PCMCI RBF & 0.8053 & 0.8752 & -0.0699 & -3.9526 & \textbf{0.000077} \\
\bottomrule
\end{longtable}}


\section{Baseline Attribution Plots}

Chapter~\ref{ch:results} reports the SHAP attributions through the relevance and consistency scores computed over them, and leaves the attributions themselves here. Table~\ref{tab:shap} gives the top-ranked variable each architecture returns for the three robotics attack classes, on the same faults as the causal attributions of Table~\ref{tab:attributions}.

% GENERATED by Thesis/tools/make_tables.py -- do not edit by hand.
% Every figure below is read from a CSV written by the pipelines.
\begin{table}[htbp]
\centering
\small
\caption[Post-hoc SHAP attribution over the neural autoencoders]{Post-hoc SHAP attribution over the neural autoencoders, on the same three robotics attack classes as the causal attributions in Table~\ref{tab:attributions}. SHAP recovers a tablet sensor for the LED attack in two of three architectures; the intrinsic causal attribution recovers one in nine of nine.}
\label{tab:shap}
\begin{tabular}{llll}
\toprule
Architecture & LedsControl & JointControl & WheelsControl \\
\midrule
LSTM & LedFaceBR5 & RShoulderPitch & WheelBSpeed \\
TCN & \textbf{TabletTouchX} & RShoulderPitch & WheelBSpeed \\
Transformer & \textbf{TabletTouchX} & RShoulderPitch & WheelBSpeed \\
\bottomrule
\end{tabular}
\end{table}


SHAP returns a different top variable depending on which architecture it is asked to explain, recovering a tablet sensor for the TCN and the Transformer on the LED fault and failing on the LSTM. The drift-based attribution returns the same answer across three algorithms with three different graphs and two different variable sets, two of which do not model the attacked sensors at all. Figures~\ref{fig:shap-lstm-leds} to~\ref{fig:shap-transformer-leds} are the three \texttt{LedsControl} panels behind that row, which is where the contrast with the causal attribution is sharpest.

\begin{figure}[htbp]
  \centering
  \includegraphics[width=0.9\textwidth]{baseline_shap_lstm_autoencoder_ledscontrol.png}
  \caption{SHAP attribution for the LSTM autoencoder on the \texttt{LedsControl} fault.}
  \label{fig:shap-lstm-leds}
\end{figure}

\begin{figure}[htbp]
  \centering
  \includegraphics[width=0.9\textwidth]{baseline_shap_tcn_autoencoder_ledscontrol.png}
  \caption{SHAP attribution for the TCN autoencoder on the \texttt{LedsControl} fault.}
  \label{fig:shap-tcn-leds}
\end{figure}

\begin{figure}[htbp]
  \centering
  \includegraphics[width=0.9\textwidth]{baseline_shap_transformer_autoencoder_ledscontrol.png}
  \caption{SHAP attribution for the Transformer autoencoder on the \texttt{LedsControl} fault.}
  \label{fig:shap-transformer-leds}
\end{figure}

The other two robotics faults follow, as Figures~\ref{fig:baseline-shap-lstm-autoencoder-jointcontrol} to~\ref{fig:baseline-shap-transformer-autoencoder-wheelscontrol}. The \texttt{JointControl} figures are the ones to read against Table~\ref{tab:deviation}: all three architectures return \texttt{RShoulderPitch}, which is the variable with the largest standardised deviation in the system and not the joint temperature the documented requirement names.

\begin{figure}[htbp]
  \centering
  \includegraphics[width=0.85\textwidth,height=0.78\textheight,keepaspectratio]{baseline_shap_lstm_autoencoder_jointcontrol.png}
  \caption[SHAP: LSTM, JointControl]{SHAP attribution for the LSTM autoencoder on the \texttt{JointControl} fault.}
  \label{fig:baseline-shap-lstm-autoencoder-jointcontrol}
\end{figure}

\begin{figure}[htbp]
  \centering
  \includegraphics[width=0.85\textwidth,height=0.78\textheight,keepaspectratio]{baseline_shap_tcn_autoencoder_jointcontrol.png}
  \caption[SHAP: TCN, JointControl]{SHAP attribution for the TCN autoencoder on the \texttt{JointControl} fault.}
  \label{fig:baseline-shap-tcn-autoencoder-jointcontrol}
\end{figure}

\begin{figure}[htbp]
  \centering
  \includegraphics[width=0.85\textwidth,height=0.78\textheight,keepaspectratio]{baseline_shap_transformer_autoencoder_jointcontrol.png}
  \caption[SHAP: Transformer, JointControl]{SHAP attribution for the Transformer autoencoder on the \texttt{JointControl} fault.}
  \label{fig:baseline-shap-transformer-autoencoder-jointcontrol}
\end{figure}

\begin{figure}[htbp]
  \centering
  \includegraphics[width=0.85\textwidth,height=0.78\textheight,keepaspectratio]{baseline_shap_lstm_autoencoder_wheelscontrol.png}
  \caption[SHAP: LSTM, WheelsControl]{SHAP attribution for the LSTM autoencoder on the \texttt{WheelsControl} fault.}
  \label{fig:baseline-shap-lstm-autoencoder-wheelscontrol}
\end{figure}

\begin{figure}[htbp]
  \centering
  \includegraphics[width=0.85\textwidth,height=0.78\textheight,keepaspectratio]{baseline_shap_tcn_autoencoder_wheelscontrol.png}
  \caption[SHAP: TCN, WheelsControl]{SHAP attribution for the TCN autoencoder on the \texttt{WheelsControl} fault.}
  \label{fig:baseline-shap-tcn-autoencoder-wheelscontrol}
\end{figure}

\begin{figure}[htbp]
  \centering
  \includegraphics[width=0.85\textwidth,height=0.78\textheight,keepaspectratio]{baseline_shap_transformer_autoencoder_wheelscontrol.png}
  \caption[SHAP: Transformer, WheelsControl]{SHAP attribution for the Transformer autoencoder on the \texttt{WheelsControl} fault.}
  \label{fig:baseline-shap-transformer-autoencoder-wheelscontrol}
\end{figure}

The SHAP attributions for the other two domains are collected here. The explainer is run on all three datasets, as Section~\ref{sec:setup-shap} describes, and the relevance and consistency figures reported in Section~\ref{sec:attribution} are computed over every anomaly class in each. One class from each of the two non-robotics domains is plotted below as a representative panel, for each of the three architectures: Figures~\ref{fig:baseline-shap-lstm-attack1} to~\ref{fig:baseline-shap-transformer-fault1}. The electrocardiogram panels rank twelve leads and the plant panels rank fifty-two process variables, so the two are not on a common axis and are not meant to be read against one another; each is read against the causal attribution for the same domain. The same caution applies between the three architectures within one domain, for the reason Section~\ref{sec:shap-normalised} sets out: the explainer returns attributions in the units of the input, so the horizontal scale of a panel is fixed by what the model reads and not by how much of the record it explains. On the electrocardiogram class plotted here the largest attribution differs by more than two orders of magnitude between the widest panel and the narrowest. What is comparable across them is the ordering of the leads, not the width of the bars.

\begin{figure}[htbp]
  \centering
  \includegraphics[width=0.85\textwidth,height=0.78\textheight,keepaspectratio]{baseline_shap_lstm_autoencoder_attack_1.png}
  \caption[SHAP: LSTM, PTB-XL]{SHAP attribution for the LSTM autoencoder on record class \texttt{Attack\_1} of the electrocardiogram domain.}
  \label{fig:baseline-shap-lstm-attack1}
\end{figure}

\begin{figure}[htbp]
  \centering
  \includegraphics[width=0.85\textwidth,height=0.78\textheight,keepaspectratio]{baseline_shap_tcn_autoencoder_attack_1.png}
  \caption[SHAP: TCN, PTB-XL]{SHAP attribution for the TCN autoencoder on record class \texttt{Attack\_1} of the electrocardiogram domain.}
  \label{fig:baseline-shap-tcn-attack1}
\end{figure}

\begin{figure}[htbp]
  \centering
  \includegraphics[width=0.85\textwidth,height=0.78\textheight,keepaspectratio]{baseline_shap_transformer_autoencoder_attack_1.png}
  \caption[SHAP: Transformer, PTB-XL]{SHAP attribution for the Transformer autoencoder on record class \texttt{Attack\_1} of the electrocardiogram domain.}
  \label{fig:baseline-shap-transformer-attack1}
\end{figure}

\begin{figure}[htbp]
  \centering
  \includegraphics[width=0.85\textwidth,height=0.78\textheight,keepaspectratio]{baseline_shap_lstm_autoencoder_fault_1.png}
  \caption[SHAP: LSTM, TEP]{SHAP attribution for the LSTM autoencoder on \texttt{Fault\_1} of the Tennessee Eastman domain.}
  \label{fig:baseline-shap-lstm-fault1}
\end{figure}

\begin{figure}[htbp]
  \centering
  \includegraphics[width=0.85\textwidth,height=0.78\textheight,keepaspectratio]{baseline_shap_tcn_autoencoder_fault_1.png}
  \caption[SHAP: TCN, TEP]{SHAP attribution for the TCN autoencoder on \texttt{Fault\_1} of the Tennessee Eastman domain.}
  \label{fig:baseline-shap-tcn-fault1}
\end{figure}

\begin{figure}[htbp]
  \centering
  \includegraphics[width=0.85\textwidth,height=0.78\textheight,keepaspectratio]{baseline_shap_transformer_autoencoder_fault_1.png}
  \caption[SHAP: Transformer, TEP]{SHAP attribution for the Transformer autoencoder on \texttt{Fault\_1} of the Tennessee Eastman domain.}
  \label{fig:baseline-shap-transformer-fault1}
\end{figure}

\begin{figure}[htbp]
  \centering
  \includegraphics[width=0.85\textwidth,height=0.78\textheight,keepaspectratio]{baseline_lstm_autoencoder_top_anomalous_variables.png}
  \caption[Reconstruction ranking: LSTM baseline]{Top anomalous variables per attack class for the LSTM autoencoder baseline, ranked by reconstruction error rather than by coefficient drift.}
  \label{fig:baseline-lstm-top}
\end{figure}


Figure~\ref{fig:baseline-lstm-top} gives the neural baseline's own ranking of anomalous variables, by reconstruction error rather than by coefficient drift.

\subsection{Normalising the SHAP Attributions}
\label{sec:shap-normalised}

The SHAP magnitudes plotted in Figures~\ref{fig:shap-lstm-leds}
to~\ref{fig:shap-transformer-leds} are small in absolute terms, and that
appearance is an artefact of the axis. GradientExplainer returns attributions in
the units of the input, so the magnitude of a panel is fixed by whatever scale
the sensors happen to be recorded in. Across the nine robot panels the largest
attribution ranges from 0.22 to 227.26, which is three orders of magnitude of
sensor scaling and not three orders of magnitude of explanation.

Normalising each panel by its own total attribution removes the scale and leaves
the share of the explanation that a variable carries, which is comparable across
architectures and faults. Figure~\ref{fig:shap-normalised} shows the LED fault on
that scale. The reference line is the 0.41\% a variable would receive if
the model spread attribution evenly over all 244 sensors. The top variable
carries between 2.3 and 68.2\% depending on the architecture, so the
attributions concentrate between 6 and 166 times more than an even spread. The
TCN is the extreme case, placing 95.9\% of its attribution on the top
three variables, while the LSTM places 6.3\% there and spreads the rest
almost flat. Read on the normalised scale the SHAP values are not small, and the
architectures differ from each other far more than the raw plots suggest.

\begin{figure}[htbp]
  \centering
  \includegraphics[width=\textwidth]{baseline_shap_normalised_ledscontrol.png}
  \caption[Normalised SHAP attribution for the \texttt{LedsControl} fault across the three neural architectures, on a common logarithmic scale]{Normalised SHAP attribution for the \texttt{LedsControl} fault across
  the three neural architectures, on a common logarithmic scale. Each bar is the
  share of that panel's total attribution. Blue marks a LED or tablet variable.
  The dashed line is the 0.41\% an evenly spreading model would assign.}
  \label{fig:shap-normalised}
\end{figure}

\section{Every Paired Attribution Comparison}
\label{app:attr-paired-full}

Chapter~\ref{ch:results} reports the paired attribution comparison at the top-three depth, which Section~\ref{sec:method-attribution} fixes as the depth every attribution result in this thesis is stated at. The other three depths are exploratory, and they are given here in full rather than summarised, because the alternative is to choose which of four depths to show after seeing what each one says.

One of them reverses the ordering. That reversal is why this table exists. At ten variables on the electrocardiogram GES leads SHAP by 9.0 points with thirteen of the fourteen classes behind it, which is a wider margin and a more lopsided split than either result reported in the chapter. It is marked degenerate and discounted for the reason Section~\ref{sec:attribution} gives: GES ranks nine of the twelve leads, so a request for the top ten returns all nine it has, its relevance is then the share of its own pool that is relevant, and that is its chance rate by construction. Its lift at that depth is 1.00 exactly. A method that returns everything is not ranking, and a comparison against it measures the pool rather than the explanation.

% GENERATED by Thesis/tools/make_tables.py -- do not edit by hand.
% Every figure below is read from a CSV written by the pipelines.
{\footnotesize
\setlength{\tabcolsep}{4pt}
\begin{longtable}{lllrrcrc}
\caption[Every paired attribution comparison, at all four reporting depths]{Every paired attribution comparison, at all four reporting depths. $\Delta$ is SHAP minus the causal algorithm, so positive favours SHAP. Sign splits the classes into those favouring SHAP, those favouring the causal side and those tied. $p$ is an exact two-sided paired sign test with ties discarded; it is Holm-adjusted within a domain at the top-three depth, which Section~\ref{sec:method-attribution} pre-specifies, and reported uncorrected at the other three, which are exploratory. A dash means the panel cannot reach the 5\% level at all, its smallest attainable $p$ being $2^{1-n}$. A dagger marks a degenerate comparison, one in which a method's selection depth reaches the pool it ranks so that it returns everything it has: GES ranks nine of the twelve electrocardiogram leads, so at ten variables it scores exactly chance by construction and the row measures the pool and not the method.}\label{tab:attr-paired-full}\\
\toprule
Dataset & Depth & Comparison & Cls. & $\Delta$ Rel.\% & Sign & $p$ & Sep. \\
\midrule
\endfirsthead
\toprule
Dataset & Depth & Comparison & Cls. & $\Delta$ Rel.\% & Sign & $p$ & Sep. \\
\midrule
\endhead
Robotics & Top-3 & SHAP $-$ CAM-UV & 3 & +11.1 & 2/0/1 & --- & --- \\
 & Top-3 & SHAP $-$ GES & 3 & +37.0 & 3/0/0 & --- & --- \\
 & Top-3 & SHAP $-$ PCMCI & 3 & +14.8 & 3/0/0 & --- & --- \\
 & Top-5 & SHAP $-$ CAM-UV & 3 & +22.2 & 2/1/0 & --- & --- \\
 & Top-5 & SHAP $-$ GES & 3 & +46.7 & 3/0/0 & --- & --- \\
 & Top-5 & SHAP $-$ PCMCI & 3 & +26.7 & 2/0/1 & --- & --- \\
 & Top-10 & SHAP $-$ CAM-UV & 3 & +17.8 & 2/1/0 & --- & --- \\
 & Top-10 & SHAP $-$ GES & 3 & +34.4 & 2/0/1 & --- & --- \\
 & Top-10 & SHAP $-$ PCMCI & 3 & +18.9 & 2/1/0 & --- & --- \\
 & $z>1$ & SHAP $-$ CAM-UV & 3 & -3.3 & 1/2/0 & --- & --- \\
 & $z>1$ & SHAP $-$ GES & 3 & +19.6 & 2/1/0 & --- & --- \\
 & $z>1$ & SHAP $-$ PCMCI & 3 & +25.0 & 2/1/0 & --- & --- \\
\midrule
Healthcare & Top-3 & SHAP $-$ CAM-UV & 14 & 0.0 & 7/6/1 & 1.0000 & no \\
 & Top-3 & SHAP $-$ GES & 14 & -6.3 & 5/7/2 & 1.0000 & no \\
 & Top-3 & SHAP $-$ PCMCI & 14 & \textbf{+11.9} & 10/1/3 & 0.0352 & yes \\
 & Top-5 & SHAP $-$ CAM-UV & 14 & +0.5 & 6/5/3 & 1.0000 & no \\
 & Top-5 & SHAP $-$ GES & 14 & -7.1 & 5/8/1 & 0.5811 & no \\
 & Top-5 & SHAP $-$ PCMCI & 14 & \textbf{+10.5} & 10/2/2 & 0.0386 & yes \\
 & Top-10 & SHAP $-$ CAM-UV & 14 & +1.0 & 6/3/5 & 0.5078 & no \\
 & Top-10 & SHAP $-$ GES$^\dagger$ & 14 & \textbf{-9.0} & 1/13/0 & 0.0018 & yes \\
 & Top-10 & SHAP $-$ PCMCI & 14 & +0.7 & 6/4/4 & 0.7539 & no \\
 & $z>1$ & SHAP $-$ CAM-UV & 14 & +2.0 & 8/5/1 & 0.5811 & no \\
 & $z>1$ & SHAP $-$ GES & 14 & -6.5 & 6/7/1 & 1.0000 & no \\
 & $z>1$ & SHAP $-$ PCMCI & 14 & +7.3 & 6/3/5 & 0.5078 & no \\
\midrule
Industrial & Top-3 & SHAP $-$ CAM-UV & 5 & -2.2 & 1/1/3 & --- & --- \\
 & Top-3 & SHAP $-$ GES & 5 & +6.7 & 1/0/4 & --- & --- \\
 & Top-3 & SHAP $-$ PCMCI & 5 & +2.2 & 2/1/2 & --- & --- \\
 & Top-5 & SHAP $-$ CAM-UV & 5 & +1.3 & 1/0/4 & --- & --- \\
 & Top-5 & SHAP $-$ GES & 5 & +6.7 & 2/0/3 & --- & --- \\
 & Top-5 & SHAP $-$ PCMCI & 5 & +4.0 & 2/0/3 & --- & --- \\
 & Top-10 & SHAP $-$ CAM-UV & 5 & 0.0 & 0/0/5 & --- & --- \\
 & Top-10 & SHAP $-$ GES & 5 & +4.0 & 2/0/3 & --- & --- \\
 & Top-10 & SHAP $-$ PCMCI & 5 & +1.3 & 2/0/3 & --- & --- \\
 & $z>1$ & SHAP $-$ CAM-UV & 5 & +3.1 & 2/2/1 & --- & --- \\
 & $z>1$ & SHAP $-$ GES & 5 & +6.4 & 3/1/1 & --- & --- \\
 & $z>1$ & SHAP $-$ PCMCI & 5 & +15.6 & 3/2/0 & --- & --- \\
\bottomrule
\end{longtable}}

